\item Equip $\mathcal Y$ and $\Delta$ with
their divisorial logarithmic structures
from $D$ and $0$. Locally, after
absorbing a unit, the map is
$t=z_1\cdots z_k$. Its monoid chart
is the diagonal
$\mathbb N\to\mathbb N^k$; it is
fine, saturated, exact and log smooth.
Thus \cite[Theorem 5.1, p.2226]{NO10}
applies in the complex analytic
category. The induced map of real
oriented logarithmic blow-ups is a
proper topological fibre bundle over
$[0,\epsilon)\times S^1$.
Fix $\theta=\arg t$. Trivializing
this bundle along the radial interval
identifies $Y_t$ homeomorphically with
the fibre $D^{\log}_\theta$ over
$(0,\theta)$. Its natural projection
\[
 p_\theta:D^{\log}_\theta\longrightarrow D
\]
is proper and surjective. Above a
point lying on $k$ local branches,
its fibre consists of the phases
$(\theta_1,\ldots,\theta_k)$ whose
product is the fixed phase $\theta$,
including the harmless phase of the
local unit. It is a torus
$(S^1)^{k-1}$, in particular connected.

\item Here is a covering-space proof that
$p_\theta$ is surjective on fundamental
groups; path lifting for the collapse
is not assumed. The normal-crossings
space $D$ is locally contractible
and has a connected universal cover
$\widetilde D$. Pull it back along
$p_\theta$. The map from the pullback
to $\widetilde D$ is still proper,
surjective and has connected fibres.
Its total space is connected: a
nontrivial separation into open and
closed sets would, by properness and
connected fibres, give such a
separation of $\widetilde D$.
On the other hand its connected
components are indexed by the cosets
of the image of
$\pi_1(D^{\log}_\theta)$ in $\pi_1(D)$.
There is therefore only one coset.
Composing $p_\theta$ with the chosen
radial homeomorphism proves the claim.

